
\chapter{Hashing}

\section{Introduction}

A recurring theme in randomized algorithms is the idea of replacing a large object by a much smaller 
\textit{fingerprint} or \textit{hash} which approximately preserves equality. Instead of comparing 
two complicated objects directly, we compare their hashes. If the hashes differ, the objects are 
certainly different. If the hashes are equal, the objects are \emph{probably} equal.\\


\noindent The power of hashing comes from the fact that computing and comparing hashes is often significantly 
faster than comparing the original objects themselves. It may also be seen as a form of compression 
that preserves some information about the original object. In data structures like the \texttt{unordered set}, 
hashing helps to achieve constant time complexity (on average) within finite memory constraints.\\


\noindent Usually hashes are many-one functions, meaning that multiple different objects can have the same hash. 
Two inputs having the same hash is called a \textit{collision}. The probability of collision depends on 
the specific hash function used and the distribution of inputs. A good hash function should minimize the 
probability of collisions for typical inputs.

\section{Universal Hash Families}

In many applications, we want a low probability of a hash collision because the performance of certain data 
structures and algorithms depends on it. Since the whole point of a hash function is to go from a large set 
(the domain $D$) to a smaller set (the range $R$), we will be interested in the case $|D| >> |R|$ where 
significant collisions are forced to occur.\\ 


\noindent Let $U$ be the domain (whose elements are called keys) and let $[m]=\{0,1,\dots,m-1\}$ denote the 
set of outputs. A family of hash functions $\mathcal H=\{h:U\to [m]\}$ is said to be 
\emph{universal} if for every pair of distinct keys $x,y\in U$, the probability of collision for a random 
$h \in \mathcal H$ is no greater than completely random assignment.

\[
\Pr_{h \in \mathcal H}[h(x)=h(y)]\le \frac1m, \;\; x \neq y
\]

The hash functions themselves are completely deterministic, so $h(x)=h(y)$ is either true or false --- its truth 
value is known if $h$ is known. So the probability is coming from the random choice of the 
hash function $h$. The key idea is that even if an adversary knows the design of the hash functions 
in the family, they cannot predict which hash function will be used so they can not force a 
collision with high probability. \\


Suppose $p$ is a prime and $U=\{0,1,\dots,p-1\}$. For any integers
$a\in\{1,\dots,p-1\}$ and $b\in\{0,\dots,p-1\}$, define
\[
h_{a,b}(x)=((ax+b)\bmod p)\bmod m
\]

The family
\[
\mathcal H=\{h_{a,b}:1\le a\le p-1,\ 0\le b<p\}
\]
forms a universal family.

\section{Strongly Universal Hash Families}

Universal hashing only bounds the probability of collisions. In many
applications, we require a stronger guarantee: not only should collisions be
unlikely, but the hash values of distinct keys should behave as independent
random variables.


A family of hash functions
\[
\mathcal H=\{h:U\to [m]\}
\]
is called \emph{strongly universal} (or \emph{2-wise independent}) if for
every pair of distinct keys $x,y\in U$ and for every pair of table locations
$i,j\in[m]$,
\[
\Pr_{h\sim\mathcal H}[h(x)=i\land h(y)=j]=\frac1{m^2}
\]


In other words, for any two distinct keys, the hash values are independent and
uniformly distributed over $[m]$.

Strong universality immediately implies universality. Indeed,
\[
\Pr[h(x)=h(y)]
=\sum_{i=0}^{m-1}\Pr[h(x)=i\land h(y)=i]
=\sum_{i=0}^{m-1}\frac1{m^2}
=\frac1m
\]

The converse, however, is not true: a universal family need not be strongly
universal.

Strongly universal hash families play an important role in randomized
algorithms, data streaming, sketching, and cryptographic protocols.

\section{\texorpdfstring{$k$}{k}-Wise Independent Hash Families}

Strong universality captures independence among pairs of keys. More generally,
we may require independence among larger collections of keys.


A family of hash functions
\[
\mathcal H=\{h:U\to [m]\}
\]
is said to be \emph{$k$-wise independent} if for every collection of distinct
keys
\[
x_1,x_2,\dots,x_k\in U
\]
and every choice of values
\[
y_1,y_2,\dots,y_k\in [m]
\]
we have
\[
\Pr_{h\sim\mathcal H}
\bigl[h(x_1)=y_1,\dots,h(x_k)=y_k\bigr]
=\frac1{m^k}
\]


Equivalently, the random variables
\[
h(x_1),h(x_2),\dots,h(x_k)
\]
are jointly independent and each is uniformly distributed over $[m]$.

The hierarchy of notions is therefore
\[
k\text{-wise independent}
\implies
(k-1)\text{-wise independent}
\implies
\cdots
\implies
2\text{-wise independent}
\implies
\text{universal}
\]

A standard construction of $k$-wise independent hash families uses
polynomials over finite fields. Let $\mathbb F_q$ be a finite field with
$q$ elements. Choose coefficients
\[
a_0,a_1,\dots,a_{k-1}\in\mathbb F_q
\]
uniformly and independently at random, and define
\[
h(x)=a_{k-1}x^{k-1}+\cdots+a_1x+a_0
\]

For any distinct points
\[
x_1,\dots,x_k\in\mathbb F_q
\]
the values
\[
h(x_1),\dots,h(x_k)
\]
are uniformly and independently distributed over $\mathbb F_q$. Hence this
construction yields a $k$-wise independent family.

Families with limited independence are extremely useful in randomized
algorithms because they often require far fewer random bits than complete
independence while still providing the guarantees needed for analysis.

\section{Polynomial Multiplication}

Suppose we are given two expressions:
\[
P(x)=\prod_{i=1}^k (a_i x+b_i), Q(x) = \sum_{j=0}^k c_j x^j
\]

We would like to verify whether
\[
P(x)\equiv Q(x)
\]

Naively expanding the product takes $\Theta(k^2)$ time, and then comparing the resulting 
polynomials takes $\Theta(k)$ time, so the total verification time is $\Theta(k^2)$.

A more efficient approach is to compute the hashes of the polynomials.
Choose three random values
\[
r_1,r_2,r_3
\]
from the domain of the polynomial.

Define the hash of a polynomial to be the tuple
\[
(P(r_1),P(r_2),P(r_3))
\]

Similarly compute
\[
(Q(r_1),Q(r_2),Q(r_3))
\]

If the hashes differ at any coordinate, then the polynomials are certainly different.

If all three evaluations agree, then with high probability the polynomials are identical.\\

The reason this works is that the nonzero polynomial $P(x)-Q(x)$ of degree at most $k$ can 
have at most $k$ roots 
in the domain from which we are choosing the random values.

If we are choosing random 32-bit integers, then the probability of a false positive 
(i.e. the polynomials are different but the hashes match) 
is at most $\left(\frac{k}{2^{32}}\right)^3$ since we are checking three independent random values, 
which is astronomically small for any reasonable $k$!\\

\noindent Each evaluation of $P$ or $Q$ at a random point takes $\Theta(k)$ time, so 
the total verification time is $\Theta(k)$. So we have achieved a significant improvement over 
the naive\footnote{There exists deterministic techniques to get this down to $\Theta(k\log k)$ time, 
but they have a large constant factor and also are very difficult to implement.} $\Theta(k^2)$ time.\\

\noindent Note that choosing 3 random values has no particular significance - we could choose any small
constant number of random values to achieve this - it is a tradeoff between time and correctness probability.

\section{Matrix Multiplication}

Suppose we are given matrices $A,B,C\in\mathbb{R}^{n\times n}$
and wish to verify whether
\[
AB=C
\]

Computing the product explicitly requires matrix multiplication. Standard matrix multiplication takes $\Theta(n^3)$ time.\\
time, while the fastest known algorithms run in roughly $\Theta(n^{2.37})$ time, though 
they are highly complicated and rarely used in practice. \\


\subsection{Freivalds' Algorithm}

Choose a random vector $X\in\{0,1,\dots,p-1\}^n$, aka of form \[\begin{bmatrix}
x_1\\
x_2\\
\vdots\\
x_n
\end{bmatrix}\]

Instead of comparing the matrices themselves, we compute $ABX$ and $CX$. Since matrix-vector 
multiplication takes only $\Theta(n^2)$ and we are doing 3 
such multiplications -- $B\cdot X$, then $A\cdot(BX)$, and $C\cdot X$.
Then we compare the two products entry by entry in $\Theta(n)$ time.
So the overall procedure is $\Theta(n^2)$ time. \\

\noindent The quantity \(MX\) may be viewed as a compact randomized hash of the matrix \(M\). If two matrices differ, 
their hashes are unlikely to coincide for a random choice of \(X\).

\noindent More generally, one may choose k independent random vectors
\[
X_1,X_2,\dots,X_k
\]
and define the hash of a matrix \(M\) as
\[
\textbf{H}(M)=(MX_1,MX_2,\dots,MX_k)
\]

Each additional random vector independently decreases the probability of error.
Here k must be $\Theta(1)$ to preserve the overall $\Theta(n^2)$ time complexity.

\subsection*{Error Probability Analysis}

Let \(M=AB-C\). We want to bound the probability that \(MX=0\) when 
\(M\) is nonzero. Choose a prime \(p\), and consider the finite field
\[
\mathbb Z_p=(\{0,1,2,\dots,p-1\},+_p,\times_p)
\]

\noindent To avoid making the entries of $M\equiv0\pmod{p}$ when it is 
nonzero over \(\mathbb Z\) (and thus avoiding making it null in $\mathbb Z_p$), we use such a $p$
which is strictly greater than every entry of $M$. We then choose the random 
vector $X\in\mathbb Z_p^n$ uniformly at random.\\

\noindent Notice that whenever \(MX=0\) over the integers, we also have \(MX=0\) over \(\mathbb Z_p\) 
since $0 \equiv 0\pmod{p}$. So the set of false positives over \(\mathbb Z\) is a subset 
of the set of false positives over \(\mathbb Z_p\). Thus the 
probability of a false positive over \(\mathbb Z\) is 
at most the probability of a false positive over \(\mathbb Z_p\).\\

\noindent Now \(\mathbb (Z_p)^n\) is a vector space\footnote{It may be helpful to brush up on 
the theory of vector spaces, but you may choose to not bother with the finer details 
and take my word for it!} of dimension \(n\) over \(\mathbb Z_p\). If $M$ is not the null matrix, 
the dimension of its nullspace\footnote{Refer \autoref{nullspace}} is at most \(n-1\).\\

\noindent A \(d\)-dimensional vector space over \(\mathbb Z_p\) contains exactly \(p^d\) vectors. So
\[
|nullspace(M)| \le p^{n-1}
\]
\vspace{0.5em}
\noindent Since \(X\) is chosen uniformly from \(\mathbb (Z_p)^n\), which contains \(p^n\) vectors,
\[
\Pr(MX=0 \mid \mathbb Z) \le \Pr(MX=0 \mid \mathbb Z_p)
\le
\frac{p^{n-1}}{p^n}
=
\frac1p
\]

A suitable choice is \hspace{0.2em} $p=2^{64}+13$ \hspace{0.2em}if we are working with 64-bit integers, 
which gives an astronomically small error probability of about $5.4\times10^{-20}$.
\section{A Matrix Problem}

\textbf{Problem Statement\footnote{This question was taken from Codeforces. The editorial indicates that my hashing solution wasn't intended, but when has that ever stopped me?}:} Given two \(n\times m\) matrices \(A\) and \(B\), determine if they are equivalent under row and column permutations. That is, can we permute the rows and columns of \(A\) to obtain \(B\)?
\\ Additional Constraints: A and B both satisfy -- the set of all entries is \{1,2, \dots, nm\}. \\

\noindent \textbf{Reduction:} First we reduce the problem to a simpler one. If we consider R(M) to be the set of all $R_i$ where each $R_i$ is the set of values in the 
$i$-th row of a matrix M, and similarly C(M) then the matrices A and B are equivalent under row and column permutations
if and only if R(A) = R(B) and C(A) = C(B).

\subsection*{Proof of Reduction}
Clearly, if A and B are equivalent under row and column permutations, then R(A) = R(B) and C(A) = C(B) 
since the row and column sets are preserved under these permutations. This proves that the set of all B such that R(B) = R(A) and C(B) = C(A)
contains all matrices equivalent to A under row and column permutations (i.e. the set of all B such that 
R(B) = R(A) and C(B) = C(A) is a superset of the set of all B equivalent to A under row and column permutations).

Now notice that the number of row and column permutations on A are $n!\cdot m!$ (including the identity permutation). 
Each of these permutations gives a different matrix since the entries are all distinct.
Also it can be shown that the number of matrices with given R(M) and C(M) is also $n!\cdot m!$ (In fact while proving 
this you will notice that the fact we are proving is quite intuitive)  

Since the two sets have the same size and one contains the other, they must be equal. 
Hence R(A) = R(B) and C(A) = C(B) implies that A and B are equivalent under row and column permutations.

(Note that we assumed all entries are distinct in the above proof)\\

\subsection{Set of Sets}

Formally, let $R_i$ be the set of values in row \(i\). 

\[
R_i = \{a_{i1}, a_{i2}, \dots, a_{im}\}
\]

The condition requires that the set of all $R_i$ of the matrix $A$ with the 
set of all $R_i$ of the matrix $B$. 

\[
\{R_1, R_2, \dots, R_n\}_A = \{R_1, R_2, \dots, R_n\}_B
\]

We require a similar condition for columns as well. This method is exact but 
computationally heavy due to the need to construct complex set-of-sets structures which 
take $\Theta(nm\log(nm))$ time  as well as a lot of space.

\begin{center}
\begin{minipage}{0.9\linewidth}
\begin{framed}
\begin{Verbatim}
set<set<int>>F;//set of sets for first matrix
for(int i=0;i<n;i++)
{
    set<int>temp(first[i].begin(),first[i].end()); //set of ith row
    F.insert(temp);
}

for(int i=0;i<n;i++)
{
    set<int>temp(second[i].begin(),second[i].end()); //set of ith row
    if (!F.count(temp)) //early exit 
    {
        cout<<"NO\n";
        return;
    }
}
\end{Verbatim}
\vspace{-4mm}
\end{framed}
\end{minipage}
\end{center}

\subsection{Hashing}

To obtain a more efficient solution, each row is hashed to a pair of values 
$(\text{sum}, \text{squaresum})$, where $\text{sum}$ and $\text{squaresum}$ denote the 
sum and sum of squares of elements in the row respectively. 

\[
h(R_i) = \left(\sum_{j=1}^{m} a_{ij}, \sum_{j=1}^{m} a_{ij}^2\right)
\]

We define a custom hash for a pair $(x,y)$ as $x \cdot 2^{32} + y$

\[
H(x,y) = x \cdot 2^{32} + y
\]

So overall the hash for the row is a single value $H(h(R_i))$.

\[
\mathbf{H}(R_i) = H(h(R_i))
\]

where $x$ and $y$ are the computed row statistics. This choice of pair hash is popular 
for 32 bit integers, as it combines the two values into a single 64-bit integer guaranteeing uniqueness. We 
insert the hashes of all rows of matrix \(A\) into an unordered set and 
then check if the hashes of rows of matrix \(B\) are present in this set. If any 
hash from \(B\) is not found in the set from \(A\), we can conclude that the matrices 
are not equivalent under row permutations.\\

\noindent Similarly for columns, we compute the column hashes and compare them in the same manner. This method has 
an expected time complexity of \(\Theta(nm)\) due to the linear time required 
to compute the hashes and the expected constant time for hash set operations.\\

\noindent This method involved hashing so the correctness of the algorithm is itself a random variable -- but 
given the constraint that the matrix entries are from \{1,2,\dots,nm\}, it seems highly 
unlikely that this will fail. In fact, it seems that this may be a unique hash in this case, but I am yet 
to find a proof for general (n,m).

\begin{center}
\begin{minipage}{0.9\linewidth}
\begin{framed}
\begin{Verbatim}
auto pair_hash=[](int x,int y){
    return ((x)<<32)+y;
};

unordered_set<long long>s;
loop(i,n)
{
    long long v1=0,v2=0;
    loop(j,m)
    {
        int x=first[i][j];
        v1+=x;
        v2+=x*x;
    }
    s.insert(pair_hash(v1,v2));
}
loop(i,n)
{
    long long v1=0,v2=0;
    loop(j,m)
    {
        int x=second[i][j];
        v1+=x;
        v2+=x*x;
    }
    if(!s.count(pair_hash(v1,v2))) //early exit
    {
        cout<<"NO\n";
        return;
    }          
}
\end{Verbatim}
\vspace{-4mm}
\end{framed}
\end{minipage}
\end{center}

\section{A Sliding Window Problem}

\textbf{The Problem} Two integer arrays $x$ and $y$, each of size $n$, are given along with an integer $k$ $(1 \le k \le n)$. For every contiguous subarray (window) of size $k$, determine whether the corresponding windows in $x$ and $y$ contain the same multiset of elements at every position.

More precisely, for every index (1-based) $i$ such that $1 \le i \le n-k+1$, we compare the multisets:
\[
\{x_i, x_{i+1}, \dots, x_{i+k-1}\} \quad \text{and} \quad \{y_i, y_{i+1}, \dots, y_{i+k-1}\}
\]
and check whether they are identical for all $i$.

\subsection*{Rolling Hash Method}

To obtain an efficient solution, we maintain two rolling hashes for each window:
\[
S = \sum a_i, \quad Q = \sum a_i^2
\]

These values can be updated in $\Theta(1)$ time when the window slides by removing the outgoing element and adding the incoming element.

For each position $i$, we compute:
\[
(S_x(i), Q_x(i)) \quad \text{and} \quad (S_y(i), Q_y(i))
\]
and check whether:
\[
(S_x(i), Q_x(i)) = (S_y(i), Q_y(i))
\]

If the condition holds for all $i$, we conclude that every corresponding window has identical multiset structure.\\

\begin{center}
\begin{minipage}{0.9\linewidth}
\begin{framed}
\begin{Verbatim}
long long sumX=0,sumY=0;
long long sumX2=0,sumY2=0;

for(int i=0;i<k;i++)
{
    sumX+=x[i];
    sumY+=y[i];
    sumX2+=x[i]*x[i];
    sumY2+=y[i]*y[i];
}

bool ok=1;
for(int i=k;i<n;i++)
{
    if(!((sumX==sumY)&&(sumX2==sumY2)))
    {
        ok=0;
        break;
    }

    sumX-=x[i-k];
    sumX+=x[i];
    sumX2-=x[i-k]*x[i-k];
    sumX2+=x[i]*x[i];

    sumY-=y[i-k];
    sumY+=y[i];
    sumY2-=y[i-k]*y[i-k];
    sumY2+=y[i]*y[i];
}
cout<<(ok?"YES":"NO");
\end{Verbatim}
\vspace{-4mm}
\end{framed}
\end{minipage}
\end{center}

This method is very easy to implement and runs in $\Theta(n)$ time, which is efficient for this problem. 
What was especially nice about this method is that we can maintain the hash values in constant time as the window slides, 
which is what I meant by using a \textbf{rolling hash}.
However, it relies on the assumption that the hash values are unique for different multisets.
We can improve this (if needed) by doing a second round of hashing with a different hash function like polynomial hashing. 

\clearpage
\section{Problems}
\begin{enumerate}

\item[\textbf{1}]
For hashes of form $H(x)=(ax+b)\bmod c$, why is it preferred to have c as a prime number? \textit{Hint: Think of 
the arithmetic progression (a very common occurence in data) with common difference d -- 
what happens when d is not coprime with c?}

\end{enumerate}